% Partie : intuitionistic-logic
% Chapitre : soundness-completeness
% Section : soundness-axd

\documentclass[../../../include/open-logic-section]{subfiles}

\begin{document}

\olfileid[fr]{int}{sc}{sax}

\olsection{Correction des \usetoken{p}{derivation} axiomatiques}

Dans un modèle relationnel, écrivons $u\leq v$ pour $Ruv$.
La vérité d'une formule se conserve quand on passe de $u$ à
$v\geq u$, par la \olref[int][sem][rel]{prop:true-monotonic}.
Pour vérifier une implication en $u$, il faut vérifier sa
conclusion en tout monde $v\geq u$ où sa prémisse est vraie.
On itère cette lecture pour les implications emboîtées.

\begin{lem}\ollabel{lem:axioms-valid}
Tous les axiomes de la déduction axiomatique intuitionniste
sont valides.
\end{lem}

\begin{proof}
Considérons les neuf schémas de la \olref[int][int][axd]{sec}.
\begin{enumerate}
\item Les deux projections de la conjonction sont valides,
  car si $!A\land !B$ est vraie en un monde, chacun de ses
  membres y est vrai. Pour l'introduction de la conjonction,
  supposons $!A$ vraie en $u$ et $!B$ vraie en $v\geq u$.
  Par monotonie, $!A$ est encore vraie en $v$, donc
  $!A\land !B$ y est vraie.
\item Les deux introductions de la disjonction sont valides
  dès que $!A$ est vraie au monde considéré. Pour l'élimination,
  prenons $u\leq v\leq w$ où sont vraies, respectivement,
  $!A\lif !C$, $!B\lif !C$ et $!A\lor !B$.
  Les deux implications sont vraies en $w$ par monotonie.
  L'un des deux disjoints est vrai en $w$ et l'implication
  correspondante donne $!C$ en $w$, par réflexivité.
\item Pour $!A\lif(!B\lif !A)$, la vérité de $!A$ se
  conserve de $u$ à tout $v\geq u$, y compris si $!B$
  est vraie en $v$. Pour le second schéma d'implication,
  prenons $u\leq v\leq w$ où sont vraies, respectivement,
  $!A\lif(!B\lif !C)$, $!A\lif !B$ et $!A$.
  Les deux premières formules sont encore vraies en $w$.
  La réflexivité donne $!B\lif !C$ puis $!B$ en $w$,
  et finalement $!C$ en $w$.
\item $\lfalse$ n'est vraie en aucun monde. L'implication
  $\lfalse\lif !A$ est donc vraie partout, par vacuité.
\end{enumerate}
\end{proof}

\begin{thm}[Correction]
  \ollabel{thm:soundness}
  Si $\Gamma \Proves !A$, alors $ \Gamma \Entails !A$.
\end{thm}

\begin{proof}
  Nous démontrons que si $\Gamma \Proves !A$, alors
  $\Gamma \Entails !A$. La preuve procède par récurrence sur le
  nombre~$n$ de !!{formula}s dans la !!{derivation} de~$!A$ à partir
  de~$\Gamma$. Nous montrons que si $!A_1$, \dots,
  $!A_n = !A$ est !!a{derivation} à partir de~$\Gamma$, alors
  $\Gamma \Entails !A_n$. Remarquons que si $!A_1$, \dots,
  $!A_n$ est !!a{derivation}, alors $!A_1$, \dots,
  $!A_k$ en est aussi une pour tout $k<n$.

  Il n'existe pas de !!{derivation} de longueur~$0$ ; pour $n=0$,
  l'énoncé est donc vrai par vacuité. Supposons-le vrai pour toutes
  les !!{derivation}s de longueur~$<n$. Nous distinguons les cas
  d'après la justification de~$!A_n$.

  \begin{enumerate}
  \item $!A_n$ est un axiome. Le \olref{lem:axioms-valid} donne
    $\Gamma \Entails !A_n$ quel que soit~$\Gamma$.
  \item $!A_n \in \Gamma$. Alors, quels que soient $\mModel{M}$ et
    $w$, si $\mSat{M}{\Gamma}[w]$, on a évidemment
    $\mSat{M}{!A_n}[w]$, c'est-à-dire $\Gamma \Entails !A$.
  \item $!A_n$ résulte de $!A_i$ et de
    $!A_j \ident !A_i \lif !A_n$ par \MP{}. Les suites $!A_1$, \dots,
    $!A_i$ et $!A_1$, \dots, $!A_j$ sont des !!{derivation}s à partir
    de~$\Gamma$. Par hypothèse de récurrence, nous avons
    $\Gamma \Entails !A_i$ et $\Gamma \Entails !A_i \lif !A_n$.

    Supposons $\mSat{M}{\Gamma}[w]$. Puisque
    $\mSat{M}{\Gamma}[w]$ et $\Gamma \Entails !A_i \lif !A_n$,
    on a $\mSat{M}{!A_i \lif !A_n}[w]$. Par définition, pour tout
    $w'$ tel que $Rww'$, si $\mSat{M}{!A_i}[w']$, alors
    $\mSat{M}{!A_n}[w']$. La relation $R$ étant réflexive,
    $w$ est l'un des $w'$ tels que $Rww'$ : ainsi,
    $\mSat{M}{!A_i}[w]$ implique $\mSat{M}{!A_n}[w]$. Comme
    $\Gamma \Entails !A_i$, nous avons $\mSat{M}{!A_i}[w]$ et donc
    $\mSat{M}{!A_n}[w]$, ce qu'il fallait démontrer.
  \end{enumerate}
\end{proof}

\end{document}
